\begin{explanationbox}
	\textbf{\textcolor{CExplanation}{一句话直觉}}

	对数中的分式 \(1+\frac{1}{n}\) 可以写成 \(\frac{n+1}{n}\)，再用对数商法则拆成相邻整数的对数差。

	\medskip
	\textbf{\textcolor{CExplanation}{核心拆解}}
	\begin{description}[style=nextline, leftmargin=2.8em, labelsep=0.5em, itemsep=0.45em]
		\item[\textbf{要点一（对数商法则）：}]
		      根据
		      \[
			      \log_a\frac{M}{N}=\log_a M-\log_a N
			      \quad (M>0,\;N>0),
		      \]
		      将左边分式看成商 \(\frac{n+1}{n}\)。

		\item[\textbf{要点二（分式恒等变形）：}]
		      \[
			      1+\frac{1}{n}=\frac{n+1}{n}.
		      \]
		      这是实现对数裂项的前提。
	\end{description}

	\medskip
	\textbf{\textcolor{CExplanation}{几何本质（数轴比例）}}

	在数轴上，从点 \(n\) 到 \(n+1\) 的长度为 \(1\)，\(\frac{1}{n}\) 可理解为这段增量相对于 \(n\) 的比例。恒等式把“比例的对数”转化为“两个位置的对数差”，体现了对数把乘除关系转化为加减关系的本质。

	\medskip
	\textbf{\textcolor{CExplanation}{代数意义（运算降级）}}

	它把对数内的除法转化为对数外的减法，把复杂的分式对数变成相邻两项之差，为后续累加消去提供基础。

	\medskip
	\textbf{\textcolor{CExplanation}{考点价值（数列裂项求和）}}
	\begin{itemize}[leftmargin=*, itemsep=0.5em, topsep=0.4em]
		\item \textbf{考法一（直接裂项）：} 通项为 \(\log_a\!\left(1+\frac{1}{k}\right)\) 时，可直接拆为 \(\log_a(k+1)-\log_a k\)。
		\item \textbf{考法二（累加消去）：} 求和时相邻项中的 \(\log_a(k+1)\) 与 \(-\log_a(k+1)\) 互相抵消，只剩首末项。
	\end{itemize}

	\medskip
	\textbf{\textcolor{CExplanation}{顿悟点}}

	看到形如 \(\log_a\!\left(1+\frac{1}{n}\right)\) 的式子，第一反应不是计算分式值，而是把它写成
	\[
		\log_a(n+1)-\log_a n.
	\]

	\medskip
	\textbf{\textcolor{CExplanation}{使用场景}}
	\begin{itemize}[leftmargin=*, itemsep=0.5em, topsep=0.4em]
		\item \textbf{场景一（数列求和）：} 求 \(\displaystyle\sum_{k=1}^{n}\log_a\!\left(1+\frac{1}{k}\right)\) 时，裂项后化简为 \(\log_a(n+1)\)。
		\item \textbf{场景二（简化表达式）：} 对数代数式化简，例如将 \(\log_a\frac{n+1}{n}\) 拆开再合并。
	\end{itemize}
\end{explanationbox}